\section{From an agent-facing toolkit to a qualified causal adapter}
\label{exec23:statspai-contract-scope}
The preceding software comparison motivates a more precise design decision.
A banking analyst may ask an agent to estimate the effect of a servicing
intervention, compare alternative estimators, and preserve a result that another
analyst can reproduce. A unified catalogue can reduce the effort of finding and
calling methods. It does not itself determine which customers, intervention,
time horizon or assignment law make the returned number relevant. StatsPAI
provides a concrete integration candidate against which to develop this
separation. The inspected cutoff source declares version 1.39.0 at commit
\texttt{35d5a4fbd05f1b90201c24d806df9b12f966e50d}, dated 7 October 2026.
This is distinct from the paper's archived version 1.23.0; an observation about
one version cannot silently be transferred to the other. The selected eleven
files include the package declaration, result-schema exporter, parity taxonomy,
MCP server, workflow guide and two test specifications.
\citep{exec23_statspai_contract,exec23_statspai_paper}

The package declares Python 3.10 or later, a substantial scientific and reporting
dependency set, optional JAX and PyTorch stacks, and separate command-line and
MCP entrypoints. Its inspected MIT notice includes an attribution condition and
warranty disclaimer. The package declaration also names third-party notices;
those notices, dependency-specific terms and dataset rights require their own
inventory before a chosen distribution. A package-level license observation
is not a conclusion about every optional component or about the separately
identified commercial downstream product. An adapter should pin the source,
selected extras and executable environment together, rather than relying on
lower bounds in a package declaration as a reproducible lock.
\citep{exec23_statspai_contract}

\subsection{Transport shape and scientific admissibility}
\label{exec23:statspai-contract-schema}
There are at least three distinct contracts. A transport contract determines
whether the response can be read. A result contract determines whether the
required fields and diagnostics for a particular method are present. A
scientific contract determines whether the estimand and uncertainty answer the
banking question under stated assumptions. The exported agent-result JSON
schema requires only a string-valued \texttt{method}. It permits method-specific
extensions, and its interval array has no length or order constraint. Its
standard error, probability and observation-count properties do not impose
nonnegativity, the unit interval or integer-valued sample size respectively.
The MCP server's compact advertised output schema is more permissive still:
it accepts any object and refers the client to a fuller result resource.
This supports a heterogeneous catalogue, but acceptance under either schema
must not be promoted into scientific admissibility.
\citep{exec23_statspai_contract}

Independent data-only checks make the boundary concrete. A payload containing
only a method name satisfies the exported schema. A second payload with a
reversed interval $[2,1]$, standard error $-1$, $p$-value $2$, error level $-1$
and observation count $0.5$ also satisfies it, as does an empty interval array.
A numeric method name fails. These are deliberately authored payloads validated
against the downloaded schema; they are not observations that an estimator
emitted such values. The check executes the project's validator and parses
upstream syntax as data; it does not import or execute StatsPAI.
\citep{exec23_statspai_contract_reference}

For a proposed average-treatment-effect adapter, the consumer contract should
require the target population, treatment contrast, units, horizon, information
cutoff and identification assumptions, together with a declared uncertainty
kind. If a finite standard error is applicable, require it to be nonnegative.
If an ordinary confidence interval is applicable, require two ordered endpoints
and an explicit coverage level. A legitimate unbounded identified set must
retain its meaning through a separate representation; turning an infinity into
an apparently missing scalar would erase information. Regression coefficient
tables, causal effects, diagnostic tests and helper responses need distinct
variants. Refusal and partial identification are substantive outcomes, not
malformed estimates to be replaced by a default value. These requirements are
our proposed consumer design, beyond the generic transport schema.

\subsection{Failures, time budgets and reproducible custody}
\label{exec23:statspai-contract-runtime}
The inspected server distinguishes JSON-RPC request failures from failures
inside a tool call. It constructs an execution-error response with an error
flag, structured error fields and a serialized text representation. The
underidentified-IV tests specify retention of diagnostic and alternative-method
information. Those tests are useful contract specifications; their presence
is not evidence that they passed in a local executable environment. Likewise,
a suggested alternative method is a proposal for analyst review, not proof
that it is identified under the original question's assumptions.
\citep{exec23_statspai_contract}

Timeout semantics particularly matter after an interrupted session. The caller
can dispatch an eligible call to a process worker, but a result that cannot be
transferred can fall back to the thread runner. The process-timeout branch
reports a killed worker. The thread branches explicitly report that a data read
or computation may remain active. Eligibility, kill behavior, handle transfer,
queue admission and orphan tracking are delegated to other modules outside this eleven-file inspection. Section~\ref{exec23:statspai-runtime-scope} examines selected additional helpers and narrows these questions; it still does not establish native lifecycle behavior. The supported conclusion is therefore
about the caller's distinct contracts, not a verified guarantee that every
call stops computing when its client stops waiting. An adapter should retain
actual execution mode and worker state, reconcile the original attempt after
reconnection, and allow a retry only after its duplication and resource
consequences are understood. This is the same distinction between an
observation deadline and terminal completion required of commercial workflow
services.
\citep{exec23_statspai_contract}

Output size has a similarly precise boundary. The server applies its budget
to the structured result and then serializes that result again in a text
content block; an image can be appended separately. The stated budget is not
a cap on the entire transmitted message. A structured JSON payload of
262,144 bytes already contributes at least 524,288 bytes when present twice,
before enclosing fields, string escaping or an image. This is an arithmetic
lower bound on the duplicated representations, not a network benchmark.
A causal adapter should measure and limit the complete response and preserve
whether risk details were shortened. The inspected caller records nonfinite
locations, truncation metadata and a capped list of runtime warnings; it calls
separate helpers for the detailed budgeting policy. A small successful message
is not sufficient if omitted diagnostics could change the decision.
\citep{exec23_statspai_contract,exec23_statspai_contract_reference}

The workflow guide recommends saving externally retrieved data and preserving
provider, indicator, query, retrieval date and transformations. It distinguishes
local-file provenance from remote data whose bytes are not hashed unless saved
locally. The caller also distinguishes a standalone replay from a replay that
depends on process-held handles or merely records a call. Thus a replay string
is not a complete reconstruction package. For a banking proposal, the reusable
artifact should contain the selected immutable input, transformation graph,
model configuration, sampling and fold choices, package environment, complete
result and diagnostics. Its permissions and confidential-data custody must be
specified separately. No sibling repository has to be modified to define this
handoff: a consumer can read a versioned artifact and test the declared result
contract in its own environment.
\citep{exec23_statspai_contract}

\subsection{Parity is conditional on the target and the statistic}
\label{exec23:statspai-contract-parity}
The taxonomy explicitly records past problems with inherited alias claims and
circular evidence. It separates estimator entries from loaders, renderers and
other infrastructure, and requires named equivalence evidence before an alias
receives another entry's grade. Its mapping labels externally aligned results
as certified, internally supported results as validated, and other states as
API-stable. These are upstream labels with specific meanings. None is an
independent certification of a bank's causal claim. The source's Python-reference
mapping contains eight entries although the adjacent comment still says two;
the literal mapping, verified without executing the module, prevents the stale
comment from narrowing the apparent cross-package exceptions. Cross-package
comparison should not be described as cross-language evidence when both
implementations are Python.
\citep{exec23_statspai_contract,exec23_statspai_contract_reference}

The inspected parity-boundary tests read committed results. For one RD fixture
they prescribe much tighter point-estimate tolerances than the allowed relative
standard-error difference of seven percent. A DML comparison specifies shared
explicit folds; a synthetic-control comparison distinguishes nonunique weight
solutions and a unique-solution counterpart. These are useful prospective
checks. They do not show that the current installation reproduced R or Stata,
or that a tight agreement on one fixture establishes coverage for another
sampling law. An adapter's evidence unit must therefore name the exact entrypoint,
input bytes, target, statistic, convention, reference version and tolerance.
The guide itself separates stored comparisons from invoking a separately
configured R or Stata runtime.
\citep{exec23_statspai_contract}

A documented non-alias illustrates why matching a method family is insufficient.
The taxonomy distinguishes cohort-size weighting from treated-observation
weighting even when both arise from related difference-in-differences estimators.
Consider our stipulated servicing example with two cohorts of ten customers
each. The first cohort contributes three post-treatment periods and has a
constant effect of USD 4 per customer-period; the second contributes one period
and has effect USD $-4$. The equal-customer average effect at a common evaluation
period is
\[
 \tau_{\mathrm{customer}}=\frac{10(4)+10(-4)}{20}=0,
\]
whereas the available treated-observation average is
\[
 \tau_{\mathrm{observation}}=\frac{10\cdot3\cdot4+10\cdot1\cdot(-4)}{40}=2.
\]
With a stipulated USD 1 per-customer decision cost, the first target gives
net USD $-1$. Subtracting that cost from the second average gives USD 1 but
mixes a different weighting target into the customer decision. Neither
arithmetic calculation repairs the mismatch. No estimator is run here, and
these dollar inputs are authored examples, not measured benefits. The adapter
must bind the aggregation and horizon before deciding whether two software
outputs disagree or simply answer different questions.
\citep{exec23_statspai_contract,exec23_statspai_contract_reference}

\subsection{Adoption decision and a bounded qualification route}
\label{exec23:statspai-contract-decision}
The preferred proposal route is a conditional isolated adapter, with StatsPAI
serving as a source-grounded reference until qualification. Its catalogue,
metadata, structured errors and replay concepts are valuable design assets.
Direct blanket adoption is deferred because the generic schema does not enforce
our scientific contract. The additional helper inspection in Section~\ref{exec23:statspai-runtime-scope} resolves part of the source uncertainty and identifies narrower lifecycle, custody and replay limitations; native guarantees remain unverified.
A selective port may be justified by a specific numerical or deployment need,
but a broad port would add solver, environment and maintenance obligations.
The rejected claim is blanket certification based on registry labels or JSON
validity; rejecting that claim does not discard the underlying package.
\citep{exec23_statspai_contract}

The executable qualification route has five bounded parts. First, inventory
third-party notices, selected dependencies and data rights, and lock one
estimator family and environment. Second, use the following same-revision helper review to specify the process, loader, cache, replay and budgeting tests, and inspect any additional selected path. Third, when a
separately permitted prototype environment exists, reproduce one named estimator
against a qualified reference with identical inputs and folds, preserving full
covariance, diagnostics and refusal behavior. Fourth, exercise stale handles,
missing files, nonfinite values, incomplete risk output, timeout, cancellation
and reconstruction after restart; scientific success must not survive an
incomplete result. Fifth, compare the adapter with a direct mature-package
wrapper on matched public or generated tasks, recording actual operational
cost separately from stipulated banking economics. Failure at one gate gives
a scoped disposition and does not block independent literature or software
review. This develops a testable product proposal while leaving native
implementation, independent acceptance and bank adoption as later evidence.
